% ============================================================================
% Section 2.9 — Financial Literacy: Budgeting, Saving, and Investing
% State-specific: Texas
% ============================================================================
\section{Financial Literacy --- Budgeting, Saving, and Investing}

\begin{stepGoal}
\begin{itemize}
    \item Create and analyze a simple monthly budget
    \item Compare saving and investing options
\end{itemize}
\end{stepGoal}

\begin{stepVocabBox}
    \vocabItem{Budget}{A plan for how you spend and save your money.}
    \vocabItem{Fixed Expense}{A cost that stays the same each month (rent, phone plan).}
    \vocabItem{Variable Expense}{A cost that changes each month (food, entertainment).}
\end{stepVocabBox}

\begin{stepsCard}[How to Create and Analyze a Budget]
    \stepItem{List your \textbf{income} (total money coming in).}
    \stepItem{Subtract \textbf{fixed expenses} and \textbf{variable expenses}.}
    \stepItem{The leftover is your \textbf{savings}. Divide by income and multiply by $100$ to find the savings rate.}
\end{stepsCard}

% --- Worked Example 1 ---
\begin{stepExample}{Monthly income: $\$1{,}200$. Fixed expenses: $\$700$. Variable expenses: $\$350$. What is the savings amount and rate?}
    \stepShow{1} Income: $\$1{,}200$.

    \stepShow{2} Subtract expenses: $1{,}200 - 700 - 350 = \$150$.

    \stepShow{3} Savings rate: $\frac{150}{1{,}200} \times 100 = 12.5\%$.
    \stepResult{$\$150$ saved, which is $12.5\%$ of income.}
\end{stepExample}

% --- Worked Example 2 ---
\begin{stepExample}{You want to save $\$1{,}200$ in $8$ months. How much do you save per month?}
    \stepShow{1} Total goal: $\$1{,}200$. Time: $8$ months.

    \stepShow{2} Divide: $1{,}200 \div 8 = \$150$.

    \stepShow{3} You need $\$150$/month. Make sure your budget leaves at least that much after expenses.
    \stepResult{Save $\$150$ per month.}
\end{stepExample}

\watchOut{Going even a little over budget each month adds up fast. Spending $\$25$ extra per month means $\$300$ less in savings per year!}

% ============================================================================
% PRACTICE
% ============================================================================
\newpage
\begin{stepPractice}{Budgeting and Saving Practice}
    \resetProblems

    \stepPracticeHeader[step1Color]{Identify Income and Expenses}
    \prob Income: $\$2{,}000$/month. Fixed: $\$1{,}100$. Variable: $\$600$. How much is left to save? \answerBlank[2cm]
    \answer{$\$300$}

    \stepPracticeHeader[step2Color]{Calculate Savings}
    \begin{multicols}{2}
    \prob You save $\$200$/month for $6$ months. How much do you have? \answerBlank[2cm]
    \answer{$\$1{,}200$}
    \prob Your food budget is $\$400$/month. You spend $\$425$. By what percent did you go over? \answerBlank[2cm]
    \answer{$6.25\%$}
    \end{multicols}

    \stepPracticeHeader[step3Color]{Compare Options}
    \prob A savings account earns $2\%$/year. An investment averages $7\%$/year but could lose money. Name one reason to choose each. \answerBlank[4cm]
    \answer{Savings: safe, low risk; Investment: higher growth potential}

    \wordProblem{Jalen earns $\$1{,}500$/month. He spends $\$850$ on fixed costs and $\$400$ on variable costs. What percent of his income does he save?}{\%}
    \answer{$\approx 16.7\%$}
\end{stepPractice}
